\chapter{Throwing Away the Target}
\label{ch:throwing_target}

\epigraph{%
  A river does not flow toward the ocean because it knows where the ocean is.
  It flows because flowing is what rivers do.
}{--- \textup{Adapted from Lao Tzu}}

\vspace{0.5cm}

\begin{laymansbox}{Setting the Stage: Motion, Not Destination}{setting_the_stage}
Classical control theory often focuses on getting a system to a specific point and keeping it there---like a thermostat holding a room at 72 degrees. However, for complex motions like a golf swing, a gymnast's routine, or a robot walking, there is no single "destination"; the motion itself is the goal. This chapter shifts the perspective from "reaching a target" to "following a trajectory." We will see how focusing on the path a system takes, rather than a single endpoint, fundamentally changes how we think about stability, optimization, and control in high-dimensional spaces. Let's explore why throwing away the target is the first step toward true mastery of motion.
\end{laymansbox}

% ──────────────────────────────────────────────
\section{The Wrong Question}
% ──────────────────────────────────────────────

Open any classical control textbook and you will find, within the first few pages, a
picture like this: a system, an error signal, a reference point, and a controller
whose entire purpose is to drive that error to zero. The implicit promise is
seductive---\emph{give me a destination and I will take you there}.

But consider the golf swing.

A skilled golfer does not steer the clubhead toward the ball the way a thermostat
drives a room toward $72^\circ$F. There is no ``setpoint'' floating in space that
the club is chasing. Instead, the golfer's body executes a \emph{motion}---a
coordinated, whole-body trajectory through a high-dimensional configuration space.
The clubhead arrives at the ball not because it was pointed at the ball, but because
the entire motion was shaped so that, somewhere along its arc, the club happens to
pass through the ball's location at enormous speed.

\begin{keyidea}{The Central Thesis}{central}
  \textbf{Control is motion, not destination.}
  The fundamental object of nonlinear control is not a setpoint, not an equilibrium,
  and not a fixed target. It is a \emph{trajectory}---a curve through state space
  that the system is asked to follow, and around which disturbances are rejected.
  The quality of control is measured not by how close the system gets to a point,
  but by how faithfully it reproduces a desired motion.
\end{keyidea}

This is not merely a philosophical reframing. It changes the mathematics. It changes
what we optimize. It changes how we define stability. And it changes what ``success''
means in a controlled system.

The purpose of this book is to develop control theory from this trajectory-first
perspective, to show that the classical setpoint framework is a special (and often
misleading) case of a much richer theory, and to build the mathematical tools needed
to design controllers for systems whose purpose is to \emph{move}.

\vspace{0.3cm}

\begin{center}
\begin{tikzpicture}[scale=1.0]
  % Classical picture
  \begin{scope}[shift={(-4.5,0)}]
    \node[font=\bfseries\color{chapblue}] at (0, 2.8) {Classical View};
    \draw[-{Stealth[length=3mm]}, thick, gray] (-1.8, 0) -- (1.8, 0)
      node[right, font=\small] {$x_1$};
    \draw[-{Stealth[length=3mm]}, thick, gray] (0, -1.8) -- (0, 1.8)
      node[above, font=\small] {$x_2$};
    % Target
    \filldraw[accentred] (0.8, 0.6) circle (3pt)
      node[right, font=\small\bfseries, accentred] {~target $\state^*$};
    % Trajectories converging
    \foreach \a/\b in {-1.3/1.4, -1.5/-0.8, 1.4/-1.2, 1.0/1.5} {
      \draw[-{Stealth[length=2mm]}, thick, chapblue!70, opacity=0.7]
        (\a, \b) .. controls (0.2, 0.1) .. (0.75, 0.57);
    }
    \node[font=\small, text width=3.5cm, align=center] at (0, -2.5)
      {All trajectories converge\\to a single point.};
  \end{scope}

  % Our picture
  \begin{scope}[shift={(4.5,0)}]
    \node[font=\bfseries\color{trajgreen!80!black}] at (0, 2.8) {Trajectory View};
    \draw[-{Stealth[length=3mm]}, thick, gray] (-1.8, 0) -- (1.8, 0)
      node[right, font=\small] {$x_1$};
    \draw[-{Stealth[length=3mm]}, thick, gray] (0, -1.8) -- (0, 1.8)
      node[above, font=\small] {$x_2$};
    % Reference trajectory
    \draw[very thick, trajgreen, decoration={markings,
        mark=at position 0.3 with {\arrow{Stealth[length=3mm]}},
        mark=at position 0.7 with {\arrow{Stealth[length=3mm]}}
      }, postaction={decorate}]
      (-1.4, -1.2) .. controls (-0.5, 0.8) and (0.5, 1.6) .. (1.5, 0.3);
    \node[trajgreen!80!black, font=\small\bfseries] at (1.2, 1.3)
      {$\traj^*(t)$};
    % Tube around trajectory
    \draw[dashed, trajgreen!50, thick]
      (-1.2, -1.5) .. controls (-0.3, 0.5) and (0.7, 1.3) .. (1.7, 0.0);
    \draw[dashed, trajgreen!50, thick]
      (-1.6, -0.9) .. controls (-0.7, 1.1) and (0.3, 1.9) .. (1.3, 0.6);
    % Nearby trajectory
    \draw[thick, accentred!70, decoration={markings,
        mark=at position 0.5 with {\arrow{Stealth[length=2mm]}}
      }, postaction={decorate}]
      (-1.3, -1.0) .. controls (-0.4, 0.9) and (0.6, 1.5) .. (1.4, 0.4);
    \node[font=\small, text width=3.5cm, align=center] at (0, -2.5)
      {System tracks a moving\\reference through state space.};
  \end{scope}

  % Dividing line
  \draw[thick, gray!40] (0, -2.0) -- (0, 2.5);
\end{tikzpicture}
\end{center}

\vspace{0.3cm}

% ──────────────────────────────────────────────
\section{Why Setpoints Are a Special Case}
% ──────────────────────────────────────────────

Let us be precise about what we mean. Consider a nonlinear system
\begin{equation}\label{eq:nonlinear_sys}
  \dot{\state} = \bm{f}(\state, \control), \qquad
  \state \in \statespace \subseteq \Reals^n, \quad
  \control \in \controlspace \subseteq \Reals^m,
\end{equation}
where $\bm{f}$ is smooth and the state space $\statespace$ may carry additional
geometric structure (it may be a manifold, as we will see in Chapter~3).

In the classical framework, the control objective is:
\begin{quote}
  \emph{Find $\control(t)$ such that $\state(t) \to \state^*$ as $t \to \infty$.}
\end{quote}
The entire apparatus of Lyapunov stability, LQR, pole placement, and PID control
is built around this question. But notice what this objective assumes:

\begin{enumerate}[leftmargin=2cm, label=\textbf{A\arabic*.}]
  \item There exists a meaningful equilibrium $\state^*$ where $\bm{f}(\state^*, \control^*) = \bm{0}$.
  \item The system's purpose is to \emph{reach and remain at} this equilibrium.
  \item Performance is measured by how quickly and smoothly the error
        $\bm{e}(t) = \state(t) - \state^*$ decays.
\end{enumerate}

Now consider the systems this book is about:

\begin{itemize}[leftmargin=1.5cm]
  \item A golfer executing a swing---the body passes through a continuum of
        configurations and never ``arrives'' anywhere.
  \item A gymnast performing a floor routine---the motion \emph{is} the performance.
  \item A biped walking---there is no single posture that constitutes ``walking'';
        walking is a \emph{limit cycle}, a periodic orbit through state space.
  \item A robotic arm tracing a weld seam---the task is defined by the path, not the endpoint.
\end{itemize}

For all of these systems, assumptions \textbf{A1}--\textbf{A3} fail. The correct
control objective is fundamentally different:

\begin{principle}{The Trajectory Tracking Objective}{tracking}
  Given a reference trajectory $\traj^*: [0, T] \to \statespace$, find a feedback
  control law $\control = \bm{k}(\state, t)$ such that
  \begin{equation}\label{eq:tracking_obj}
    \norm{\state(t) - \traj^*(t)} \leq \varepsilon(t)
    \qquad \text{for all } t \in [0, T],
  \end{equation}
  where $\varepsilon(t)$ is a prescribed tracking tolerance (the ``tube width''),
  and the system is robust to disturbances $\bm{d}(t)$ satisfying
  $\norm{\bm{d}(t)} \leq \delta$.
\end{principle}

A setpoint is simply the degenerate case where $\traj^*(t) = \state^*$ for all $t$
and $T \to \infty$. The trajectory-tracking framework \emph{contains} the setpoint
framework, but not vice versa.

\begin{definition}{Trajectory Tube}{tube}
  Given a reference trajectory $\traj^*: [0,T] \to \statespace$ and a tube radius
  function $\varepsilon: [0,T] \to \Reals_{>0}$, the \textbf{trajectory tube} is the
  time-varying set
  \begin{equation}
    \tube(t) = \bigl\{\, \state \in \statespace : \norm{\state - \traj^*(t)} \leq \varepsilon(t) \,\bigr\}.
  \end{equation}
  A controller achieves \textbf{tube-invariance} if every solution starting in
  $\tube(0)$ remains in $\tube(t)$ for all $t \in [0, T]$.
\end{definition}

This definition is the trajectory-centric replacement for Lyapunov stability. Instead
of asking ``does the system converge to a point?'' we ask ``does the system stay
inside a moving tube?'' The tube can be tight where precision matters (e.g., at
ball impact in the golf swing) and loose where it does not (e.g., during the
backswing).

% ──────────────────────────────────────────────
\section{Anatomy of a Motion: The Golf Swing as Prototype}
% ──────────────────────────────────────────────

We will use the golf swing throughout this book as a running example---not because
golf is inherently more interesting than other motions, but because it compresses
an extraordinary number of control-theoretic challenges into roughly 1.2~seconds:

\begin{enumerate}[label=(\roman*)]
  \item \textbf{High dimensionality.} The human body has over 200 degrees of freedom.
        Even a simplified musculoskeletal model of the golf swing uses 20--40
        generalized coordinates.
  \item \textbf{Extreme nonlinearity.} Joint torques couple through the full rigid-body
        dynamics, including Coriolis and centrifugal terms that dominate at high
        angular velocities.
  \item \textbf{Underactuation.} Once the club is released in the downswing, the
        wrist hinge ``freewheels''---the golfer cannot directly control the club angle.
        The clubhead speed at impact comes from \emph{passive dynamics}, not
        active torque.
  \item \textbf{Time-criticality with flexible timing.} The swing must produce the
        right state (clubhead position, velocity, and orientation) at impact,
        but the \emph{exact} time of impact is not prescribed. The golfer
        is free to take the backswing slowly or quickly.
  \item \textbf{Repeatability over optimality.} A touring professional does not
        execute the theoretically optimal swing. They execute a swing they
        can \emph{repeat} 300 times in a tournament, under pressure, with
        minimal variance.
\end{enumerate}

\begin{example}{From Setpoint Thinking to Trajectory Thinking}{golf_shift}
  Suppose we model a simplified golf swing with generalized coordinates
  $\bm{q} = (\theta_{\text{torso}},\, \theta_{\text{shoulder}},\,
  \theta_{\text{arm}},\, \theta_{\text{wrist}})^\top$ and their velocities
  $\dot{\bm{q}}$, giving state $\state = (\bm{q}, \dot{\bm{q}}) \in \Reals^8$.

  \textbf{Setpoint approach} (wrong): Define a target state $\state^* =
  (\bm{q}_{\text{impact}}, \dot{\bm{q}}_{\text{impact}})$ and design a controller
  to drive $\state(t) \to \state^*$. This fails immediately: the system should not
  \emph{stop} at $\state^*$. Impact is a waypoint, not a destination. The club must
  be accelerating \emph{through} the ball, not decelerating toward it.

  \textbf{Trajectory approach} (right): Define a reference motion
  $\traj^*(s): [0, 1] \to \Reals^8$ parameterized by a phase variable $s$ (not
  clock time $t$). The controller's job is to ensure that the system traverses
  $\traj^*$ with the right sequencing and coordination, with a tight tube near the
  impact phase ($s \approx 0.7$) and a loose tube during setup ($s \approx 0$).
\end{example}

Point~(v) deserves special emphasis. Classical optimal control asks: ``What is
the minimum-time or minimum-energy trajectory?'' But this is the wrong question
for biological and athletic systems. The right question is:

\begin{principle}{The Repeatability Principle}{repeatability}
  The best trajectory is not the one that maximizes performance on a single trial,
  but the one that maximizes \emph{expected performance over many trials} in the
  presence of noise, fatigue, and uncertainty. Formally, we seek
  \begin{equation}\label{eq:repeatability}
    \traj^* = \arg\min_{\traj} \;
    \mathbb{E}\!\left[\, \int_0^T
      \ell\bigl(\state(t), \control(t)\bigr) \dt \,\right]
    + \lambda \, \mathrm{Var}\!\left[\, J(\traj, \bm{w}) \,\right],
  \end{equation}
  where $J$ is the performance cost, $\bm{w}$ represents stochastic disturbances
  (motor noise, perceptual error), and $\lambda > 0$ penalizes variance.
\end{principle}

This is why elite golfers do not all swing the same way. Each golfer finds a
trajectory that is locally optimal \emph{for their body and their noise characteristics}.
Robustness is not an afterthought bolted onto an optimal controller---it is the
primary design objective.

% ──────────────────────────────────────────────
\section{The Mathematical Landscape}
% ──────────────────────────────────────────────

To make the trajectory-first perspective rigorous, we need mathematical machinery
that most introductory control courses skip entirely. This section provides a roadmap
of the key concepts we will develop in subsequent chapters.

\subsection{Phase Variables and Time Reparameterization}

A crucial distinction separates \emph{what} the system does from \emph{when} it
does it. Consider two executions of the same golf swing: one with a slow backswing
and one with a fast backswing. The \emph{spatial path} through configuration space
may be identical, but the \emph{time parameterization} differs.

We formalize this by introducing a \textbf{phase variable} $s(t) \in [0, 1]$ that
monotonically increases from $0$ (start of motion) to $1$ (end of motion), with
dynamics
\begin{equation}\label{eq:phase}
  \dot{s} = \alpha(s, \state),
\end{equation}
where $\alpha > 0$ is the \emph{phase velocity}. The reference trajectory is defined
as a function of $s$, not $t$:
\begin{equation}
  \traj^* = \traj^*(s), \qquad s \in [0, 1].
\end{equation}
This decouples the \emph{shape} of the motion from its \emph{temporal execution}.
The controller has two separate tasks:
\begin{enumerate}[label=(\alph*)]
  \item \textbf{Path tracking}: keep $\state(t)$ close to the reference path
        $\traj^*(s(t))$, and
  \item \textbf{Timing control}: regulate $\dot{s}$ to achieve the desired
        speed profile along the path.
\end{enumerate}
This decomposition is impossible in the setpoint framework, where time and space
are fused together.

\subsection{Transverse Dynamics and Orbital Stability}

When we shift from point stability to trajectory stability, the relevant notion of
``closeness'' changes. We do not care about the distance from $\state(t)$ to
$\traj^*(t)$ at the same clock time---we care about the distance from $\state(t)$
to the \emph{nearest point on the trajectory}.

\begin{definition}{Transverse Distance}{transverse}
  Let $\traj^*: [0,1] \to \statespace$ be a smooth reference trajectory. The
  \textbf{transverse distance} from a state $\state$ to $\traj^*$ is
  \begin{equation}
    d_\perp(\state, \traj^*) = \min_{s \in [0,1]} \norm{\state - \traj^*(s)}.
  \end{equation}
  The minimizing argument $s^*(\state) = \arg\min_s \norm{\state - \traj^*(s)}$
  is the \textbf{projection phase}, giving the ``closest point on the reference.''
\end{definition}

This leads to a decomposition of the state into two components:
\begin{itemize}[leftmargin=1.5cm]
  \item The \textbf{tangential component}: where the system is \emph{along} the
        trajectory (measured by $s$).
  \item The \textbf{transverse component}: how far the system is \emph{from} the
        trajectory (measured by $d_\perp$).
\end{itemize}

\begin{center}
\begin{tikzpicture}[scale=1.2]
  % Reference trajectory
  \draw[very thick, trajgreen, decoration={markings,
      mark=at position 0.15 with {\arrow{Stealth[length=3mm]}},
      mark=at position 0.55 with {\arrow{Stealth[length=3mm]}},
      mark=at position 0.85 with {\arrow{Stealth[length=3mm]}}
    }, postaction={decorate}]
    (-3, -0.5) .. controls (-1, 1.5) and (1, 1.5) .. (3, -0.3);
  \node[trajgreen!70!black, font=\small\bfseries] at (3.3, 0.1) {$\traj^*(s)$};

  % Point on trajectory
  \coordinate (P) at (0.3, 1.45);
  \filldraw[trajgreen!70!black] (P) circle (2pt);
  \node[above left, font=\small, trajgreen!70!black] at (P) {$\traj^*(s^*)$};

  % Actual state
  \coordinate (X) at (0.9, 0.2);
  \filldraw[accentred] (X) circle (2pt);
  \node[below right, font=\small\bfseries, accentred] at (X) {$\state(t)$};

  % Transverse distance
  \draw[thick, accentred, dashed, {Stealth[length=2.5mm]}-{Stealth[length=2.5mm]}]
    (P) -- (X) node[midway, right, font=\small] {$d_\perp$};

  % Tangent vector
  \draw[-{Stealth[length=3mm]}, thick, chapblue]
    (P) -- ++(0.9, -0.15) node[above, font=\small] {tangential};

  % Phase annotation
  \draw[thin, gray, |->] (-3, -1.3) -- (3, -1.3)
    node[right, font=\small] {$s$};
  \node[font=\small, gray] at (-3, -1.6) {$0$};
  \node[font=\small, gray] at (3, -1.6) {$1$};
  \draw[thin, gray, dashed] (0.3, -1.3) -- (P);
  \filldraw[gray] (0.3, -1.3) circle (1.5pt);
  \node[font=\small, gray, below] at (0.3, -1.3) {$s^*$};
\end{tikzpicture}
\end{center}

\vspace{0.2cm}

A system is \textbf{orbitally stable} around $\traj^*$ if the transverse distance
$d_\perp(\state(t), \traj^*)$ remains small (or decays) even if the tangential
phase $s(t)$ drifts relative to the nominal timing. This is exactly the right notion
for the golf swing: we want the motion to follow the right \emph{path} even if the
golfer is slightly faster or slower than planned.

\subsection{Funnel Control and Time-Varying Tubes}

Combining the trajectory tube (Definition~\ref{def:tube}) with the
transverse/tangential decomposition leads to what we call \textbf{funnel control}:
the design of controllers that keep the system inside a prescribed time-varying tube
around the reference.

The tube width $\varepsilon(s)$ is a design parameter. For the golf swing:
\begin{equation}\label{eq:funnel_profile}
  \varepsilon(s) = \begin{cases}
    \varepsilon_{\text{loose}} & s \in [0,\, 0.3] \quad \text{(address \& backswing)}, \\[4pt]
    \varepsilon_{\text{loose}} - (\varepsilon_{\text{loose}} - \varepsilon_{\text{tight}})
      \dfrac{s - 0.3}{0.4} & s \in [0.3,\, 0.7] \quad \text{(transition \& downswing)}, \\[6pt]
    \varepsilon_{\text{tight}} & s \in [0.7,\, 0.8] \quad \text{(impact zone)}, \\[4pt]
    \text{relaxed} & s \in [0.8,\, 1.0] \quad \text{(follow-through)}.
  \end{cases}
\end{equation}

\begin{center}
\begin{tikzpicture}[xscale=8, yscale=2.5]
  % Axes
  \draw[-{Stealth[length=3mm]}, thick] (0, 0) -- (1.08, 0)
    node[right, font=\small] {phase $s$};
  \draw[-{Stealth[length=3mm]}, thick] (0, 0) -- (0, 1.3)
    node[above, font=\small] {tube width $\varepsilon(s)$};

  % Funnel profile
  \draw[very thick, accentred]
    (0, 1.0) -- (0.3, 1.0)
    -- (0.7, 0.15)
    -- (0.8, 0.15)
    (0.8, 0.15) .. controls (0.85, 0.3) and (0.9, 0.7) .. (1.0, 0.9);

  % Phase labels
  \node[font=\scriptsize, below, text width=1.3cm, align=center] at (0.15, -0.05)
    {backswing\\(loose)};
  \node[font=\scriptsize, below, text width=1.5cm, align=center] at (0.5, -0.05)
    {downswing\\(narrowing)};
  \node[font=\scriptsize, below, text width=1.2cm, align=center, accentred] at (0.75, -0.05)
    {\textbf{impact}\\(tight)};
  \node[font=\scriptsize, below, text width=1.5cm, align=center] at (0.92, -0.05)
    {follow-through\\(relaxed)};

  % Dashed guidelines
  \foreach \x in {0.3, 0.7, 0.8} {
    \draw[thin, gray, dashed] (\x, 0) -- (\x, 1.1);
  }
\end{tikzpicture}
\end{center}

This funnel profile encodes the physical insight that precision matters most at
impact and least during setup. The controller ``spends'' its control authority
where it matters, rather than trying to maintain uniform precision everywhere---a
lesson that classical setpoint control cannot teach.

% ──────────────────────────────────────────────
\section{What Changes When Control Is Motion}
% ──────────────────────────────────────────────

Adopting the trajectory-first perspective changes nearly every aspect of control
system design. Here we preview the major consequences, each of which will be
developed in detail in subsequent chapters.

\subsection{Stability Is Not Convergence---It Is Confinement}

In the setpoint framework, stability means convergence: the state approaches the
equilibrium and stays near it. In the trajectory framework, stability means
\emph{confinement}: the state remains inside the moving tube.

\begin{definition}{Trajectory Stability}{traj_stability}
  A reference trajectory $\traj^*$ is \textbf{stable} under the closed-loop
  dynamics $\dot{\state} = \bm{f}(\state, \bm{k}(\state, t))$ if for every
  $\varepsilon > 0$ there exists $\delta > 0$ such that
  \begin{equation}
    d_\perp(\state(0), \traj^*) < \delta \implies
    d_\perp(\state(t), \traj^*) < \varepsilon \quad \forall\, t \in [0, T].
  \end{equation}
  It is \textbf{asymptotically trajectory-stable} if, additionally,
  $d_\perp(\state(t), \traj^*) \to 0$ as the system evolves.
\end{definition}

Note the subtle but critical difference: we use the transverse distance $d_\perp$,
not the time-indexed error $\norm{\state(t) - \traj^*(t)}$. This allows the system
to be ahead of or behind schedule without being considered ``unstable.''

\subsection{The Cost Function Is a Functional, Not a Function}

In setpoint control, the cost is typically a function of the error $\bm{e}(t)$:
\begin{equation}
  J_{\text{setpoint}} = \int_0^\infty \bigl(\bm{e}^\top \bm{Q}\, \bm{e}
  + \control^\top \bm{R}\, \control\bigr) \dt.
\end{equation}
In trajectory control, the cost is a \textbf{functional} that operates on the
entire motion:
\begin{equation}\label{eq:traj_cost}
  J_{\text{traj}}[\state(\cdot), \control(\cdot)] =
  \underbrace{\int_0^T w_\perp(s)\, d_\perp^2\bigl(\state(t), \traj^*\bigr) \dt}_{%
    \text{tracking fidelity (transverse)}}
  + \underbrace{\int_0^T w_\parallel(s)\, (\dot{s} - \dot{s}_{\text{ref}})^2 \dt}_{%
    \text{timing fidelity (tangential)}}
  + \underbrace{\int_0^T \control^\top \bm{R}(s)\, \control \dt}_{%
    \text{effort}},
\end{equation}
where the weighting matrices $w_\perp(s)$, $w_\parallel(s)$, and $\bm{R}(s)$ are
\emph{phase-dependent}. This is what allows us to demand high precision at impact
and low effort during the backswing.

\subsection{Feedforward Is Not Optional---It Is Primary}

In the setpoint framework, feedforward is a ``nice to have'' that reduces transient
error. In the trajectory framework, feedforward is the \emph{primary} control
action. The reference trajectory $\traj^*(t)$ implicitly defines a feedforward
control signal
\begin{equation}
  \control_{\text{ff}}(t) = \bm{f}^{-1}\!\left(\dot{\traj}^*(t),\, \traj^*(t)\right),
\end{equation}
where $\bm{f}^{-1}$ denotes the inverse dynamics. The feedback controller handles
only the \emph{deviations} from this feedforward signal:
\begin{equation}\label{eq:ff_fb_decomp}
  \control(t) = \underbrace{\control_{\text{ff}}(t)}_{\text{executes the motion}}
  + \underbrace{\control_{\text{fb}}(\state(t), t)}_{\text{corrects deviations}}.
\end{equation}
For the golf swing, the feedforward component accounts for approximately 90\% of
the total joint torques. The feedback controller provides the remaining 10\% of
correction. A control design that treats everything as feedback is asking the
feedback loop to do ten times more work than necessary.

\subsection{Underactuation Creates Geometry, Not Just Constraints}

Many moving systems are underactuated: they have fewer control inputs than degrees
of freedom. In the setpoint framework, underactuation is an obstacle---it means we
cannot independently control all states. In the trajectory framework, underactuation
creates \emph{geometric structure} that we can exploit.

The set of states reachable from a given configuration through passive dynamics
forms a submanifold of the state space. A well-designed trajectory \emph{rides along}
this manifold, using the passive dynamics to generate motion that active torques
alone could not produce. This is how the golf swing works: the wrist hinge releases
passively, and the centrifugal acceleration of the downswing \emph{automatically}
squares the clubface through the double pendulum effect.

\begin{proposition}
  For a mechanical system with $n$ degrees of freedom and $m < n$ actuators, the
  set of dynamically feasible trajectories lies on a submanifold of dimension
  at most $2m$ in the $2n$-dimensional state space. Effective trajectory design
  requires understanding the geometry of this submanifold.
\end{proposition}

We will make this precise in Chapter~5 using the language of constraint manifolds
and their null-space structures.

% ──────────────────────────────────────────────
\section{A Roadmap for the Book}
% ──────────────────────────────────────────────

The remainder of this book develops the trajectory-first theory in a sequence of
increasingly sophisticated layers:

\begin{description}[leftmargin=1.5cm, labelwidth=1.2cm, font=\bfseries\color{chapblue}]
  \item[Ch.~2] \textbf{Curves in State Space.}
    The geometry of trajectories: parameterization, arc length, curvature, and torsion
    in high-dimensional state spaces. Frenet--Serret frames for nonlinear systems.

  \item[Ch.~3] \textbf{Configuration Manifolds.}
    Why state space is often not $\Reals^n$: Lie groups, joint spaces, and the
    geometry of articulated bodies. The golf swing lives on $SO(3)^k$, not $\Reals^k$.

  \item[Ch.~4] \textbf{Orbital Stability and Transverse Linearization.}
    Rigorous definitions of trajectory stability. Linearization transverse to the
    orbit. Floquet theory for periodic motions. Moving Poincar\'e sections.

  \item[Ch.~5] \textbf{Underactuation and Passive Dynamics.}
    Constraint Jacobians, null-space control, and the exploitation of passive
    mechanical coupling. The double pendulum and the physics of the wrist release.

  \item[Ch.~6] \textbf{Trajectory Optimization.}
    Direct collocation, shooting methods, and pseudospectral methods for finding
    optimal reference trajectories. The interplay between optimality and
    repeatability.

  \item[Ch.~7] \textbf{Funnel Synthesis.}
    Designing time-varying tubes with guaranteed invariance. Sums-of-squares
    programming for funnel computation. Robust funnels under model uncertainty.

  \item[Ch.~8] \textbf{Phase-Variable Control}~\cite{WesterveltGrizzle2007}.
    Decoupling path shape from timing. Virtual constraints and hybrid zero dynamics
    for walking and running. Central pattern generators as phase oscillators.

  \item[Ch.~9] \textbf{Stochastic Trajectories and Motor Variability.}
    Signal-dependent noise in biological actuators. The minimum-variance theory of human movement~\cite{TodorovJordan2002}. Why optimal trajectories are smooth~\cite{FlashHogan1985}.

  \item[Ch.~10] \textbf{Learning to Move.}
    Iterative learning control, policy search, and trajectory libraries. How to
    improve a motion over repeated trials without losing stability.

  \item[Ch.~11] \textbf{Case Study: The Complete Golf Swing.}
    Full application of the entire theory to a 15-DOF musculoskeletal model of the
    golf swing, from trajectory optimization through funnel synthesis to
    stochastic robustness analysis.
\end{description}

% ──────────────────────────────────────────────
\section{A New Language for Control}
% ──────────────────────────────────────────────

Before proceeding, let us establish the conceptual vocabulary that will replace
the classical language throughout this book:

\begin{center}
\renewcommand{\arraystretch}{1.4}
\begin{tabular}{@{} l @{\qquad$\longrightarrow$\qquad} l @{}}
  \textbf{\color{gray}Classical term} & \textbf{\color{chapblue}Trajectory term} \\
  \hline
  Setpoint / reference & Reference trajectory $\traj^*(s)$ \\
  Error $\bm{e}(t) = \state - \state^*$ & Transverse deviation $d_\perp(\state, \traj^*)$ \\
  Lyapunov stability & Orbital / trajectory stability \\
  Convergence to equilibrium & Confinement to tube \\
  Regulation & Path tracking + timing control \\
  Step response & Trajectory tracking response \\
  Steady-state error & Persistent transverse offset \\
  Gain margin & Funnel robustness margin \\
  Pole placement & Transverse eigenvalue assignment \\
\end{tabular}
\end{center}

\vspace{0.5cm}

% ──────────────────────────────────────────────
\section{Exercises}
% ──────────────────────────────────────────────

\begin{enumerate}[label=\textbf{1.\arabic*.}, leftmargin=1.5cm]

\item \textbf{(Conceptual.)} Consider a simple pendulum $\ddot{\theta} + \sin\theta = u$.
  \begin{enumerate}[label=(\alph*)]
    \item Identify the equilibria and classify their stability.
    \item Now consider the task of making the pendulum execute a full revolution
          (pumping up from rest to continuous spinning). Explain why this task
          \emph{cannot} be formulated as setpoint regulation.
    \item Sketch the reference trajectory in the $(\theta, \dot{\theta})$ phase
          plane for one complete revolution. What is the natural phase variable?
  \end{enumerate}

\item \textbf{(Mathematical.)} For the system $\dot{x}_1 = x_2$,\;
  $\dot{x}_2 = -x_1 + u$, consider the reference trajectory
  $\traj^*(t) = (\cos t,\, -\sin t)^\top$ (a unit circle in state space).
  \begin{enumerate}[label=(\alph*)]
    \item Verify that $u_{\text{ff}}(t) = 0$ produces exactly this trajectory.
    \item Compute the transverse linearization about $\traj^*$.
    \item Design a feedback law $u_{\text{fb}}$ that achieves asymptotic orbital
          stability.
  \end{enumerate}

\item \textbf{(Design.)} A simplified 2-DOF golf swing model has coordinates
  $\bm{q} = (\theta_{\text{arm}}, \theta_{\text{club}})$ with dynamics
  \[
    \bm{M}(\bm{q})\ddot{\bm{q}} + \bm{C}(\bm{q}, \dot{\bm{q}})\dot{\bm{q}}
    + \bm{g}(\bm{q}) = \bm{B}\,\control,
    \qquad \bm{B} = \begin{pmatrix} 1 \\ 0 \end{pmatrix}.
  \]
  This system is underactuated: only $\theta_{\text{arm}}$ is directly actuated.
  \begin{enumerate}[label=(\alph*)]
    \item Explain qualitatively how the club can accelerate through passive dynamics
          even though no torque is applied at the wrist.
    \item Define a phase variable $s$ and sketch a funnel profile $\varepsilon(s)$
          appropriate for this system.
    \item Discuss why a setpoint controller applied at impact ($s = 0.7$) would
          produce the wrong motion.
  \end{enumerate}

\item \textbf{(Philosophical.)} A thermostat, an autopilot, and a dancer's body
  are all ``controlled systems.'' Classify each according to whether its primary
  control objective is (i)~setpoint regulation, (ii)~trajectory tracking, or
  (iii)~something else entirely. Justify your classification.

\end{enumerate}

\vspace{1cm}
\begin{center}
  \rule{0.5\textwidth}{0.4pt}\\[0.5cm]
  {\small\itshape
  The river does not arrive at the ocean and stop.\\
  It becomes the ocean.\\
  The swing does not arrive at impact and stop.\\
  It becomes the follow-through.\\[0.3cm]
  Control is motion.}
\end{center}

